\documentclass[conference]{IEEEtran}
\usepackage[T1]{fontenc}
\usepackage[utf8]{inputenc}
\IfFileExists{brazil.ldf}{
  \usepackage[brazil]{babel}
}{
  \usepackage{babel}
  \babelprovide[main,import=pt-BR]{brazil}
}
\usepackage{microtype}
\usepackage{amsmath}
\usepackage{booktabs}
\usepackage{tabularx}
\usepackage{array}
\usepackage{url}
\usepackage{tcolorbox}
\newcolumntype{Y}{>{\raggedright\arraybackslash}X}
\newcommand{\pendencia}[1]{\textbf{[PREENCHER: #1]}}


\title{Reúso Exato e Evolução de Agent Skills no GitHub}
\author{\IEEEauthorblockN{Jemison dos Santos}
\IEEEauthorblockA{Universidade Estadual de Maringá -- PR, Brasil\\
\texttt{jemison321@gmail.com}}}


\begin{document}
\maketitle

\begin{abstract}
\textbf{Contexto:} Agent Skills são artefatos versionáveis cujo compartilhamento não implica rastreabilidade. \textbf{Objetivo:} Caracterizar reúso exato, evolução posterior à âncora e candidatos a mudanças potencialmente comportamentais. \textbf{Métodos:} A analise do GitSkills foi feita com estatística descritiva, regressão logística e PPML condicional; em amostra histórica condicionada de casos completos, foram comparados \textit{patches} e classificações estrita/ampla. \textbf{Resultados:} Entre 1.586.188 conteúdos, 16,02\% ocorreram em múltiplos repositórios. As 1.151.636 presenças adicionais representam 42,06\% dos 2.737.824 pares conteúdo/repositório. Os modelos utilizaram 1.584.628 conteúdos. Dos 55 repositórios e 122 ocorrências selecionados, restaram 53 repositórios e 82 ocorrências completas; 80 tiveram âncora resolvida e seis divergiram. Encontrou-se somente um candidato de propagação de confiança C, insuficiente para demonstrar transmissão; a taxa não pôde ser estimada. Em 75 comparações posteriores à âncora, de 18 repositórios, as frequências de candidatos comportamentais foram sensíveis às regras. Nenhum modelo multivariado da RQ3 foi viável. \textbf{Conclusões:} O reúso é minoritário entre conteúdos, mas substancial nas presenças. Os resultados longitudinais são exploratórios e condicionados à amostra; identidade, similaridade, ancestralidade e propagação exigem evidências distintas.
\end{abstract}

\begin{IEEEkeywords}
Agent Skills, reúso exato, mineração de repositórios, evolução de software
\end{IEEEkeywords}

\section{Introdução}
Agent Skills organizam instruções e recursos para agentes em uma pasta centrada em \texttt{SKILL.md}, opcionalmente acompanhada por scripts e referências~\cite{agentskills-spec}. Tratá-las como artefatos versionáveis implica considerar interface, dependências e manutenção~\cite{authoring-skills}. Entretanto, conteúdo compartilhado sem rastreabilidade dificulta distinguir atualização coordenada, adaptação independente e influência de fontes comuns não observadas.

Reúso exato significa presença do mesmo conteúdo, identificado por hash, em repositórios distintos. Similaridade admite diferenças; evolução descreve alterações de uma ocorrência; ancestralidade requer evidência de descendência; propagação exige evidência de transmissão de alterações. Igualdade de conteúdo ou patch não estabelece essas últimas relações. Analogamente, comandos e URLs no texto não comprovam execução.

O GitSkills permite investigar reúso e manutenção~\cite{gitskills}. A lacuna abordada é articular alcance transversal, evolução recuperável e força da evidência inter-repositórios. O objetivo é caracterizar esses fenômenos preservando suas diferentes exigências de evidência:
\begin{itemize}
\item \textbf{RQ1:} Como o reúso exato de Agent Skills está distribuído e quais características textuais, estruturais e contextuais estão associadas à sua ocorrência e intensidade?
\item \textbf{RQ2:} Entre ocorrências com histórico completo e âncora temporal resolvida, quais padrões de estabilidade e divergência são observados e que evidências existem de propagação de alterações entre repositórios?
\item \textbf{RQ3:} Quais mudanças potencialmente comportamentais são identificadas após a âncora e quão sensíveis suas frequências são às definições estrita e ampla?
\end{itemize}
Contribuímos com uma caracterização do reúso exato e suas associações, uma reconstrução exploratória de evolução posterior à âncora, uma busca de candidatos de propagação com confiança explícita e uma avaliação da sensibilidade de regras e dos limites da inferência multivariada.


\section{Trabalhos Relacionados}

O GitSkills reúne arquivos \texttt{SKILL.md}, suas ocorrências, hashes, metadados e parte do histórico de alterações, oferecendo uma base para investigar o reúso e a manutenção de Agent Skills~\cite{gitskills}. Essa possibilidade se aproxima da perspectiva de Destefanis~\cite{authoring-skills}, que propõe tratar skills como artefatos de software, considerando seu versionamento, suas interfaces e dependências. Sob essa perspectiva, compreender o reúso envolve examinar tanto a presença de um conteúdo em diferentes repositórios quanto as alterações observadas ao longo de sua manutenção.

A literatura sobre clones de código oferece uma aproximação para essa análise. Krinke~\cite{krinke2007} e Bettenburg et al.~\cite{bettenburg2012} investigam mudanças consistentes e inconsistentes em clones, destacando a importância do contexto histórico. No caso das skills, porém, essa aproximação exige cautela: os artefatos combinam linguagem natural, comandos e recursos relacionados, e a igualdade entre conteúdos ou patches não é suficiente para estabelecer ancestralidade ou transmissão de alterações.

A presença desses diferentes elementos também motiva o exame de mudanças potencialmente comportamentais. A estrutura prevista na especificação de Agent Skills~\cite{agentskills-spec} permite considerar o documento e os arquivos relacionados, mas identificar comandos, URLs ou referências a scripts não demonstra que tenham sido executados. Por isso, as classificações automatizadas deste estudo são tratadas como candidatos cuja frequência depende das regras adotadas.

Este trabalho articula essas perspectivas ao analisar a distribuição do reúso exato, a evolução posterior à âncora de conteúdo e a sensibilidade das classificações automatizadas. A análise distingue presença compartilhada, evolução e propagação, explicitando os denominadores e os limites das evidências disponíveis para cada questão.


\section{Método}

Realizamos um estudo observacional sobre o snapshot do GitSkills. O processo foi dividido em uma etapa transversal, destinada a caracterizar o reúso exato (RQ1), e uma etapa longitudinal exploratória, destinada a reconstruir a evolução posterior das ocorrências (RQs 2 e 3). Esta última foi condicionada à acessibilidade dos repositórios e à recuperação completa de seus históricos Git.

\subsection{Preparação dos dados e análise transversal}

Inicialmente, selecionamos registros com \texttt{dedup\_primary=1}, nome exato \texttt{SKILL.md}, conteúdo recuperado, hash validado e \textit{front matter} válido. Definimos \textbf{conteúdo} como o texto identificado por \texttt{file\_sha} e \textbf{ocorrência} como a presença desse conteúdo em um caminho de um repositório. Assim, hashes iguais representaram reúso exato, independentemente do caminho do arquivo. Os metadados contextuais foram os do repositório representante da deduplicação, sem interpretá-lo como autor ou origem do conteúdo.

Para cada conteúdo $h$, calculamos o alcance $R(h)$ como o número de repositórios distintos em que ele ocorreu. A ocorrência de reúso foi definida por $I\{R(h)>1\}$; sua intensidade, por $R(h)-1$, somente entre conteúdos reutilizados. A participação das presenças adicionais foi calculada por $\sum_h[R(h)-1]/\sum_hR(h)$, tendo pares conteúdo--repositório como denominador. Extraímos como covariáveis tamanho do texto e da descrição, legibilidade Flesch, URLs, blocos de código, cabeçalhos, listas, scripts, referências, arquivos irmãos, localização e composição truncada, além de linguagem, estrelas, idade e inatividade do repositório representante. Excluímos a variável de fork e registros com covariáveis ausentes.

Descrevemos a distribuição do reúso por frequências, percentis e coeficiente de Gini. Em seguida, ajustamos um modelo em duas partes: regressão logística para estimar a associação das covariáveis com a ocorrência de reúso e Poisson de pseudomáxima verossimilhança (PPML), com ligação logarítmica, para estimar a média condicional da intensidade positiva. Nas variáveis numéricas assimétricas, aplicamos \texttt{log1p}, winsorização nos percentis 0,5 e 99,5 e padronização. Reportamos \textit{odds ratios} na logística e razões de médias condicionais no PPML, por desvio-padrão da variável preparada, com IC de 95\%, erros-padrão agrupados pelo representante e correção Benjamini--Hochberg separada por componente ($q<0{,}05$). Avaliamos a sensibilidade mediante exclusão de megarepositórios, limitação da cauda positiva e comparação entre a amostra de teste e a população modelada.

\subsection{Recuperação histórica e análise longitudinal}

Para a RQ2, construímos o quadro amostral com 603 repositórios, pertencentes a 55 conteúdos-âncora do piloto, que permaneceram acessíveis em duas sondagens consecutivas. Em cada um de 11 estratos primários mutuamente exclusivos, selecionamos cinco repositórios sem reposição por ordenação pseudoaleatória, usando a semente \texttt{20261026}. Repositórios com recuperação incompleta foram excluídos integralmente, sem substituição; os pesos preservaram as probabilidades de seleção no quadro condicionado.

Para cada ocorrência, percorremos o histórico e fixamos como \textbf{âncora} o primeiro timestamp observado cujo blob coincidiu com o hash do snapshot. Comparamos cada commit posterior com cada um de seus pais, inclusive separando os pais de merges, e excluímos do cálculo de divergência ocorrências sem âncora resolvida. Classificamos como estáveis as ocorrências sem alteração de conteúdo na janela recuperada e distinguimos alterações de bytes de movimentos que preservaram o conteúdo.

Dos diffs posteriores à âncora, produzimos assinaturas \texttt{exact-normalized-patch-v1.0}, preservando operação, direção e ordem das linhas após normalização de Unicode e espaços. Excluímos movimentos puros e patches sem linhas alteradas. Um candidato a propagação exigiu a mesma família-âncora, divergência posterior, repositórios distintos e ordem temporal estrita; repetições foram consolidadas na primeira ocorrência do patch em cada repositório. Classificamos a evidência como A (proveniência explícita ou ancestralidade Git compartilhada), B (mesma transição entre blobs não vazios), C (patch exato menos específico, incluindo adição ou remoção) ou D (evidência insuficiente). Estimamos a divergência não ponderada com IC de Wilson e a ponderada com bootstrap estratificado por repositório; a taxa de propagação somente seria estimada quando o contexto anterior do patch fosse aplicável às linhagens comparadas.

\subsection{Classificação das mudanças e validação}

Na RQ3, aplicamos regras multirrótulo aos diffs para identificar candidatos a mudanças em comandos, acesso externo, scripts e executáveis. A definição estrita incluiu comandos, operações de rede, invocações explícitas, modo executável e \textit{shebang}; a ampla acrescentou URLs, caminhos e linhas compatíveis com comandos. Comparamos suas frequências, verificamos se o conjunto estrito estava contido no amplo e calculamos a similaridade de Jaccard.

Organizamos uma linha por comparação, com 20 desfechos e covariáveis anteriores ao evento. Para cada desfecho, exigimos ao menos 50 eventos, dez positivos, cinco repositórios positivos, cinco negativos e 30 clusters para modelagem multivariada. Quando o critério não foi atendido, aplicamos apenas contrastes descritivos ponderados, com bootstrap estratificado por repositório e 2.000 réplicas, sem testes confirmatórios.

Por fim, validamos o pipeline com testes automatizados, \textit{fixtures}, invariantes e reexecuções determinísticas. Registramos versões, parâmetros e checksums das entradas e saídas. Essa validação verificou a reprodução das regras, mas não sua validade semântica, pois não houve avaliação humana independente nem execução controlada.


\section{Resultados}
\subsection{RQ1: distribuição e associações do reúso exato}
O reúso exato foi observado em uma parcela minoritária dos conteúdos analisados. Dos 1.586.188 conteúdos distintos, 1.332.110 (83,98\%) apareceram em apenas um repositório, enquanto 254.078 (16,02\%) ocorreram em múltiplos repositórios. 

Apesar dessa menor frequência entre conteúdos, as 1.151.636 presenças adicionais corresponderam a 42,06\% dos 2.737.824 pares conteúdo/repositório observados. A distribuição do alcance apresentou assimetria: a mediana foi de um repositório por conteúdo e a média de 1,73, enquanto o máximo chegou a 937. Os percentis mostram que o alcance permaneceu reduzido para a maior parte dos conteúdos: P50 e P75 foram iguais a um, P90 correspondeu a dois repositórios e P95 a três. Na extremidade superior, P99, P99,5 e P99,9 alcançaram, respectivamente, 15, 29 e 65 repositórios. Apenas 466 conteúdos ultrapassaram 100 repositórios.

O coeficiente de Gini de 0,3978 complementa essa caracterização da desigualdade do alcance, sem que a cauda longa observada estabeleça uma distribuição de lei de potência. A participação dos conteúdos de maior alcance foi mais expressiva nas presenças adicionais do que no total de pares conteúdo/repositório. Os dez primeiros responderam por 0,5877\% das presenças adicionais e 0,2476\% dos pares. Para os 100 primeiros, essas participações foram de 2,7300\% e 1,1520\%; para os 1.000 primeiros, de 10,8565\% e 4,6032\%, respectivamente.

Na análise das características associadas ao reúso, os modelos retiveram 1.584.628 conteúdos, vinculados a 218.408 repositórios representantes, após a exclusão de 1.560 registros com dados ausentes (aproximadamente 0,10\%). O PPML considerou apenas o subconjunto com reúso positivo.

Os resultados diferiram entre ocorrência e intensidade como vemos na Tabela~\ref{tab:efeitos}. Após a correção BH, cinco efeitos principais atenderam ao critério de significância na ocorrência e seis na intensidade. A quantidade de cabeçalhos apresentou associação positiva com ambos os desfechos, enquanto o tamanho da descrição e o índice Flesch apresentaram associações negativas. Arquivos irmãos e inatividade atenderam ao critério somente na ocorrência; URLs, blocos de código e estrelas, somente na intensidade. A presença de scripts e referências não apresentou associação estatisticamente significativa após a correção em nenhum dos componentes.

\begin{table}[h!]
\caption{Efeitos principais com $q_{\mathrm{BH}}<0{,}05$}
\label{tab:efeitos}
\centering\footnotesize
\begin{tabular}{@{}lccc@{}}
\toprule
Característica & Razão & IC de 95\% & $q_{\mathrm{BH}}$ \\
\midrule
\multicolumn{4}{@{}l}{\textit{Ocorrência: OR}} \\
Descrição & 0,924 & 0,869--0,982 & 0,0482 \\
Flesch & 0,778 & 0,721--0,839 & $<0{,}0001$ \\
Cabeçalhos & 1,241 & 1,145--1,345 & $<0{,}0001$ \\
Arquivos irmãos & 1,121 & 1,029--1,220 & 0,0458 \\
Inatividade & 1,223 & 1,071--1,395 & 0,0211 \\
\midrule
\multicolumn{4}{@{}l}{\textit{Intensidade positiva: razão PPML}} \\
Descrição & 0,867 & 0,814--0,924 & 0,0001 \\
Flesch & 0,903 & 0,866--0,942 & $<0{,}0001$ \\
URLs & 1,075 & 1,048--1,103 & $<0{,}0001$ \\
Blocos de código & 0,920 & 0,877--0,964 & 0,0023 \\
Cabeçalhos & 1,106 & 1,043--1,172 & 0,0026 \\
Estrelas & 0,838 & 0,762--0,922 & 0,0013 \\
\bottomrule
\end{tabular}
\par\smallskip
\begin{minipage}{\columnwidth}\footnotesize
Efeitos por desvio-padrão na escala preparada; erros agrupados pelo representante. OR descreve chances, não variação em pontos percentuais. BH separado por componente.
\end{minipage}
\end{table}

Os erros agrupados foram, na mediana, 10,95 vezes os HC1. Na parte positiva, média=4,53, variância=154,30 e dispersão de Pearson=16,47. Excluir 16 megarepositórios com mais de 5.000 conteúdos (183.623 registros; 11,59\%) preservou as direções dos 28 efeitos principais, com mudanças de significância limítrofes e maior instabilidade nas linguagens. Limitar em 70 os 1.267 casos acima de P99,5 preservou direção e classificação nominal ($p<0{,}05$) dos 37 coeficientes PPML, sem demonstrar invariância BH; diferença média absoluta=0,0308 e dispersão=9,23. Entre teste de 100.000 registros e população, 71/74 coeficientes HC1 preservaram direção, com diferença média absoluta=0,0684.

\begin{tcolorbox}[
    colback=gray!15,
    colframe=gray!50,
    boxrule=0.4pt,
    arc=1mm,
    left=2mm,
    right=2mm,
    top=1mm,
    bottom=1mm
]
\textbf{Resposta da RQ1: } O reúso exato ocorreu em 16,02\% dos conteúdos, mas suas presenças adicionais representaram 42,06\% dos pares conteúdo--repositório. Mais cabeçalhos, descrições menores e menor índice Flesch estiveram associados à ocorrência e à intensidade do reúso; as demais características apresentaram associações específicas a cada desfecho. Esses resultados são associativos e não indicam causalidade.
\end{tcolorbox}


\subsection{RQ2: estabilidade, divergência e propagação}
O fluxo foi 55 repositórios/122 ocorrências selecionados, 53/82 completos, 80 âncoras resolvidas e seis ocorrências divergentes; separadamente, observaram-se 75 comparações posteriores à âncora. Essas unidades não são intercambiáveis. A divergência não ponderada foi 7,5\% (IC Wilson de 95\%: 3,5--15,4\%); a ponderada condicional, aproximadamente 1,0\% (bootstrap: 0,3--1,9\%). Pesos desiguais e concentração explicam diferenças; nenhuma estimativa representa toda a população.

As comparações incluem 27 adições, 24 modificações, 19 movimentos e cinco remoções: 36,0\%, 32,0\%, 25,3\% e 6,7\% não ponderados; 48,7\%, 12,7\%, 37,3\% e 1,3\% ponderados. Os limites bootstrap por operação não foram fornecidos; o manuscrito registra amplitudes elevadas especialmente para adições e movimentos.

Dos 173 eventos históricos, 151 produziram assinaturas elegíveis; 19 grupos apareceram em múltiplos repositórios considerando todo o histórico. Após filtros temporais, de família e colapso por repositório, restou um candidato, antes artificialmente repetido em quatro pares. Trata-se de remoções compatíveis separadas por aproximadamente 77,12 dias, classificadas como confiança C. Remover o mesmo conteúdo produz naturalmente patches iguais, de baixa especificidade; não houve relações A/B. Precedência observada não atribui origem.


\begin{tcolorbox}[
    colback=gray!15,
    colframe=gray!50,
    boxrule=0.4pt,
    arc=1mm,
    left=2mm,
    right=2mm,
    top=1mm,
    bottom=1mm
]
\textbf{Resposta da RQ2: } Predominou a estabilidade: 74 das 80 ocorrências com âncora resolvida não divergiram, enquanto seis apresentaram divergência. Foi identificado apenas um candidato de propagação de baixa confiança (nível C), sem evidência de transmissão entre os repositórios. Portanto, não houve evidência forte de propagação, e sua taxa permaneceu não estimável.
\end{tcolorbox}

\subsection{RQ3: candidatos e sensibilidade das regras}
As 75 comparações abrangem 25 ocorrências, 18 repositórios e 17 famílias-âncora; 48 possuem blob pai. A Tabela~\ref{tab:categorias} reúne as categorias com contagens disponíveis. Comandos adicionados passam de 14,7\% a 53,3\%; referências a scripts adicionadas, de 1,3\% a 28,0\%; e acesso externo adicionado, de 4,0\% a 12,0\%, entre estrito e amplo. Jaccard baixo indica muitos positivos exclusivos da definição ampla, sem determinar qual cenário é semanticamente correto.

\begin{table}[h!]
\caption{Candidatos multirrótulo em 75 comparações}
\label{tab:categorias}
\centering\footnotesize
\begin{tabularx}{\columnwidth}{@{}Yrrr@{}}
\toprule
Categoria & Estrito & Amplo & Jaccard \\
\midrule
Comando adicionado & 11 & 40 & 0,275 \\
Comando removido & 5 & 16 & 0,313 \\
Acesso externo adicionado & 3 & 9 & 0,333 \\
Acesso externo removido & 1 & 1 & 1,000 \\
Referência a script adicionada & 1 & 21 & 0,048 \\
Referência a script removida & 4 & 11 & 0,364 \\
Executável relacionado adicionado & 1 & 1 & 1,000 \\
Executável relacionado movido/copiado & 0 & 1 & 0,000 \\
\bottomrule
\end{tabularx}
\par\smallskip
\begin{minipage}{\columnwidth}\footnotesize
Contagens não aditivas; rótulos não comprovam execução. Oito pares disponíveis, não a totalidade dos 20 desfechos do gate. Repositórios positivos por rótulo não foram informados.
\end{minipage}
\end{table}

No amplo, 16/40 eventos de comando adicionado também adicionam referência a script; todos os nove de acesso externo adicionado também recebem comando adicionado; sete dos 21 de referência a script também recebem acesso externo. Uma linha pode acionar vários rótulos, sem mecanismos independentes.

Um repositório concentra 29/75 eventos. Nenhum dos 20 desfechos passou pelo gate: 18 clusters são menos que 30, independentemente dos demais requisitos. Os 280 contrastes entre 14 exposições e 20 desfechos são descritivos, não ajustados, sem $p$ ou FDR. Idade, reúso, conteúdo do pai e características concorrentes exibem contrastes sensíveis a pesos, concentração e regras; churn e merge também ampliam a oportunidade de detectar candidatos.

\begin{tcolorbox}[
    colback=gray!15,
    colframe=gray!50,
    boxrule=0.4pt,
    arc=1mm,
    left=2mm,
    right=2mm,
    top=1mm,
    bottom=1mm
]
\textbf{Resposta da RQ3: } As frequências de mudanças potencialmente comportamentais foram sensíveis às definições estrita e ampla, com mais candidatos identificados pela definição ampla. Nenhum desfecho apresentou suporte amostral para modelagem multivariada. Assim, os resultados são descritivos e não comprovam execução nem associações independentes.
\end{tcolorbox}

\section{Discussão}
Na RQ1, alcance e presença adicional descrevem composição, sem revelar circulação temporal. Uma hipótese não testada é que cabeçalhos acompanhem documentação reutilizável; outra é que reflitam convenções de agregadores. Flesch negativo não recomenda reduzir legibilidade, nem estrelas negativas indicam que menor popularidade favoreça adoção, pois o contexto é o do representante.

Na RQ2, estabilidade pode refletir manutenção de referências estáveis ou observação limitada: são hipóteses não distinguidas pelos dados. A atenção ao histórico converge com estudos de clones~\cite{krinke2007,bettenburg2012}, sem transferir seus resultados para skills. O candidato C admite decisões independentes ou fonte comum; contexto desconhecido impede tratar toda ausência de patch como oportunidade de transmissão.

Na RQ3, uma hipótese é que a definição ampla misture referências documentais e instruções operacionais, enquanto a estrita omita formas relevantes. Sem validação, nenhuma é verdade semântica. Coocorrência e concentração reforçam que 75 eventos não equivalem a 75 observações independentes.

Para autores e mantenedores, os limites motivam registrar origem declarada, versão, licença, hash, dependências e histórico de atualização, coerentemente com o tratamento de skills como software~\cite{authoring-skills}; sua eficácia não foi testada. Para pesquisadores, unidades, perdas, pesos, taxonomias e resultados não estimáveis devem ser explícitos. Futuras coletas devem ampliar clusters e janelas, verificar contexto aplicável a cada patch, incorporar proveniência e objetos Git compartilhados, validar classificações independentemente e reaplicar o gate antes de regressões. Similaridade ampliada pode recuperar candidatos, mas não substituir evidência histórica.

\section{Ameaças à Validade}
As definições adotadas limitam a interpretação dos resultados. O hash identifica conteúdo exato, sem abranger equivalência semântica, e o Flesch tem limitações para textos técnicos e multilíngues. Da mesma forma, a âncora não informa quando houve incorporação, e as classificações automatizadas não comprovam execução. Sem validação independente, suas taxas de erro permanecem desconhecidas. A seleção também restringe seu alcance: snapshot, filtros, acessibilidade, perda de 40 ocorrências, duas âncoras não resolvidas e janelas desiguais condicionam os resultados. Os pesos não corrigem essas perdas nem autorizam generalização ao GitHub. Essas restrições exigem cautela na inferência. Com poucos casos e clusters, os contrastes da RQ3 permanecem descritivos. As reexecuções verificam repetibilidade, mas não validade semântica; sua reprodução futura ainda depende da recuperabilidade dos repositórios e da resolução das pendências de logs e manifestos.
Ferramentas de inteligência artificial generativa foram utilizadas como apoio à revisão textual e à organização do manuscrito. Esse apoio não constitui validação das análises, e a responsabilidade pela conferência das afirmações, dos valores e das referências permanece com o autor.

\section{Conclusão}
Na RQ1, reúso exato é minoritário entre conteúdos, mas substancial nas presenças, com associações distintas de ocorrência e intensidade. Na RQ2, 06/80 ocorrências divergiram após a âncora; houve somente um candidato C, sem evidência forte de propagação ou taxa estimável. Na RQ3, classificações foram sensíveis às regras e faltou suporte amostral a todos os modelos multivariados. A contribuição é explicitar resultados e limites: presença compartilhada, similaridade, ancestralidade e propagação não são equivalentes.

Os artefatos estão disponíveis em \url{https://github.com/Jemisson/gitskills-reuse-analysis}; banco obtido separadamente conforme os autores do GitSkills~\cite{gitskills}.

\section{Use of Generative Artificial Intelligence}
Generative AI tools were used to review scripts and correct spelling in the manuscript. All content was critically reviewed by the author, who assumes full responsibility for the originality and accuracy of the final text, in accordance with CNPq Ordinance No. 2,664/2026.


\begin{thebibliography}{00}
\bibitem{gitskills}
G. Destefanis, D. Graziotin, M. Vaccargiu e M. Ortu, ``GitSkills: A Dataset of Agent Skills on GitHub,'' em \emph{Proc. 24th Int. Conf. on Mining Software Repositories (MSR)}, 2027, a publicar. arXiv:2608.10906.
\bibitem{agentskills-spec}
Agent Skills, ``Agent Skills Specification,'' especificação aberta, Apache-2.0. [Online]. Disponível: \url{https://agentskills.io}. Acesso em: 26 out. 2026.
\bibitem{authoring-skills}
G. Destefanis, ``Authoring Agent Skills: A Software-Engineering Approach,'' arXiv:2607.25032, 2026.
\bibitem{krinke2007}
J. Krinke, ``A Study of Consistent and Inconsistent Changes to Code Clones,'' em \emph{Proc. 14th Working Conf. on Reverse Engineering (WCRE)}, 2007, pp. 170--178.
\bibitem{bettenburg2012}
N. Bettenburg, W. Shang, W. M. Ibrahim, B. Adams, Y. Zou e A. E. Hassan, ``An Empirical Study on Inconsistent Changes to Code Clones at the Release Level,'' \emph{Science of Computer Programming}, vol. 77, pp. 760--776, 2012, doi: 10.1016/j.scico.2010.11.010.
\end{thebibliography}
\end{document}
